% Main-text section: Schubert's printed class.  Input file; no preamble.
% Macros introduced here (prefix pc): \pcXzero = X_0, \pcLog = l, \pcw = omega, \pcNorm = the printed norm.
\newcommand{\pcXzero}{X_0}
\newcommand{\pcLog}{\ell}
\newcommand{\pcw}{\omega}
\newcommand{\pcNorm}[1]{\norm{#1}_{\mathrm P}}

\section{Schubert's printed class}\label{sec:printed}

Throughout, $\Gamma$ is a fixed Fuchsian group, $M=\Hb/\Gamma$, $x_0\in\Hb$ is the base point of the norms below, and
$d=d(x_0,\cdot)$. Schubert's printed norm is \cite[(1.15)]{S}
\[
  \pcNorm{a}:=\norm{\ue^{\beta\langle d\rangle}\ue^{\alpha\sqrt{-\Delta}}a}_{L^2(\Hb)},\qquad
  \langle d\rangle=(1+d^2)^{1/2}-1,\qquad \beta>\tfrac12 .
\]

\subsection{The critical weight}\label{subsec:pc-weights}
Let $\pcXzero:=\cosh d$, $\psi:=\log\pcXzero\ge0$, $\pcLog:=\log(2\pcXzero)/\log2\ge1$ and $\Lambda:=\sqrt{-\Delta}$. Fix $c':=\tfrac34$ (any exponent in
$(\tfrac12,1)$ would do) and put
\[
  \pcw:=\pcXzero^{1/2}\pcLog^{\,c'} .
\]
This is the periodization weight: $\sum_\gamma\pcw(\gamma z)^{-2}$ is bounded for every Fuchsian group (Proposition~\ref{os:p1}); for $\pcXzero^{1/2}$ alone
the sum diverges when $M$ has finite area. Since
$\psi\le d\le\langle d\rangle+1$, $\pcw\le C_\beta\ue^{\beta\langle d\rangle}$ for every $\beta>\tfrac12$. So an amplitude of the printed class is
$a=\ue^{-\alpha\Lambda}f$ with $\norm{\pcw f}\le C_\beta\pcNorm a$: its data $f$ lie in the weighted space of the critical weight, and no weight above the
critical one is available. In the swapped class the mollifier part of the error is small because one vector is small without a weight and bounded with
the weight $\pcXzero^\beta$, $\beta>\tfrac12$. Here the same interpolation is carried out for the linear maps that produce these vectors from the data,
between $L^2(\Hb)$ and $L^2(\pcXzero^{2}d\nu)$ (Lemma~\ref{os:transfer}), after the data have been written as $a=\ue^{-\sigma\Lambda_m}\tilde f$ with an even
smoothing operator and $\norm{\pcw\tilde f}\le C\pcNorm a$ (Proposition~\ref{os:data}).

\subsection{The theorem}\label{subsec:pc-thm}
\begin{theorem}[printed class]\label{pc:main}
Let $\Gamma$ be a Fuchsian group (cusps allowed), $\beta>\tfrac12$ and $\kappa>0$. There are constants $C,\delta>0$, depending on $\Gamma$, $\beta$,
$\kappa$ and the base point $x_0$,
such that for all $0<\hbar\le1$, all $\alpha\in(0,\tfrac12]$ with $\alpha^3\ge\kappa\hbar$, all $K,N\in\mathbb N_0$, all $b\in\partial\Hb$, all $a$ with
$\pcNorm{a}<\infty$, and all $0\le t\le\delta(\alpha/\hbar)^{1/2}$,
\[
  \norm{\Delta^N\big[\cU(t)\big(a\,\ue^{i\varphi_b/\hbar}\big)_\Gamma-u^{(K)}(t)_\Gamma\big]}_{L^2(M)}
  \le C^{N+K+1}N^{2N}K!\,(\alpha\hbar)^{-2N}\Big(\frac{\hbar t}{\alpha^2}\Big)^{K+1}\pcNorm{a},
\]
with $N^{2N}:=1$ at $N=0$.
\end{theorem}

\begin{remark}
On the window of \cite[(1.31)]{S}, $t\le\delta\alpha^{3/2}\hbar^{-1/2}$, the bound of Theorem~\ref{pc:main} implies \cite[(1.30)]{S} with the printed
norm, because $(\alpha\hbar)^{4}\hbar t/\alpha^2\le1$ there. The swapped-class version is Theorem~\ref{thm:dd-appx} of Appendix~\ref{sec:dd}; the proof
below uses its frame and its lemmas. The proof gives $\delta=\min(\tfrac1{64},c_A/6)/\sqrt6$, with $c_A$ the constant of Proposition~\ref{prop:a0}
for $\tau_0=1$; this $\delta$ does not depend on $\Gamma$, $\beta$ or $\kappa$.
\end{remark}

\begin{remark}[the two classes, and the range of $\alpha$]\label{pc:classes}
Let $\beta>\tfrac12$ and $0\le\alpha'<\alpha<\pi/2$, and put $h:=\ue^{-\beta\langle d\rangle}\mathbf 1_{B(x_0,1)}$ and $a:=\ue^{-\alpha\Lambda}h$, so that
$\pcNorm a=\norm{\mathbf 1_{B(x_0,1)}}<\infty$ at $(\alpha,\beta)$. For $0<\tau<\pi/2$ the function $\ue^{-\tau\Lambda}h$ is the convolution of $h$ with the
radial kernel $k_\tau$ of Lemma~\ref{os:poisson}, and $d(z,w)\in[d(z)-1,d(z)+1]$ for $w\in B(x_0,1)$. So by \eqref{os:lb}, for $d(z)\ge2$,
\[
  (\ue^{-\tau\Lambda}h)(z)\ge c\,(d(z)+1)^{-3/2}\ue^{-d(z)},
\]
with $c>0$ depending on $\tau$ and $\beta$. Since $\langle d\rangle\ge d-1$, for every $\beta'>\tfrac12$ the square of
$\ue^{\beta'\langle d\rangle}\ue^{-\tau\Lambda}h$ is at least $c'\ue^{(2\beta'-2)d}(d+1)^{-3}$ where $d\ge2$, which is not integrable against the area
element $\sinh\rho\,d\rho\,d\theta$; so this function is not in $L^2(\Hb)$. Two consequences follow.
With $\tau=\alpha$: the amplitude $a$ of the printed class lies in no swapped class with an exponent $\beta'>\tfrac12$, because
$\norm{\ue^{\alpha''\Lambda}\ue^{\beta'\langle d\rangle}a}\ge\norm{\ue^{\beta'\langle d\rangle}a}$ for $\alpha''\ge0$.
With $\tau=\alpha-\alpha'$ and $\beta'=\beta$: the printed norm of $a$ at $(\alpha',\beta)$ is infinite, since
$\ue^{\alpha'\Lambda}a=\ue^{-(\alpha-\alpha')\Lambda}h$. So the printed norms are not monotone in $\alpha$, and Theorem~\ref{pc:main}, which is stated for
$\alpha\le\tfrac12$, does not cover larger radii by a comparison of norms; in the swapped class it does (Section~\ref{sec:repair}).
\end{remark}

\begin{corollary}[Theorem 1.3 for the printed class]\label{pc:thm13}
Under the hypotheses of Theorem~\ref{pc:main}, with $K=N=0$,
\[
  \norm{\cU(t)\big(a\,\ue^{i\varphi_b/\hbar}\big)_\Gamma-u^{(0)}(t)_\Gamma}_{L^2(M)}\le C\,\frac{\hbar t}{\alpha^2}\,\pcNorm{a},
  \qquad 0\le t\le\delta(\alpha/\hbar)^{1/2}.
\]
\end{corollary}

\subsection{Proof of Theorem~\ref{pc:main}}\label{subsec:pc-proof}
We use the notation and the statements of Appendices~\ref{app:kl}, \ref{app:ws}, \ref{sec:dd} and~\ref{app:o1st}; all lemmas are used with $\tau_0=1$.

\medskip\noindent\emph{Step 0 (reductions and parameters).} By Lemma~\ref{lem:base} we may take $x_0=i$ and $b=\infty$: the group enters below only through
the constant $C(\Gamma,c')$ of Proposition~\ref{os:p1}, which is invariant under the conjugation of Lemma~\ref{lem:base}. Since
$a\in\operatorname{Dom}(\ue^{\alpha\Lambda})$, $a$ is $C^\infty$ by Lemma~\ref{lem:lift}. Put
\[
  m:=3,\qquad\sigma:=\alpha/6,\qquad E:=\ue^{-\sigma\Lambda_m},\qquad\varepsilon:=(\sigma\hbar)^{1/2}/2,\qquad x:=(\sigma/\hbar)^{1/2},
\]
so that $\sigma\le\tfrac1{12}$, $\varepsilon\le\tfrac12$, $\varepsilon/\hbar=x/2$, $\sigma/\varepsilon=2x$ and $\sigma^3\ge\kappa'\hbar$ with $\kappa':=\kappa/216$. For
$\mathbf s\in\Sigma_{K+1}(t)$ let $B'_{\mathbf s}:=\Delta_\varphi(s_1)^N\Delta(s_1)\cdots\Delta(s_{K+1})$, so that $b'(\mathbf s)=B'_{\mathbf s}a$ in (D4) of
Proposition~\ref{prop:dd-id}. For $n\ge0$ put
\[
  F_n:=C_5^{N+K+n+1}(2N)!\,(n+2K+2)!\,\sigma^{-(n+2K+2)}(\sigma\hbar)^{-2N},
\]
with a constant $C_5\ge1$ fixed in Step~2(a).

\medskip\noindent\emph{Step 1 (the data).} By Proposition~\ref{os:data} with $c=c'$ (here $\sigma m=\alpha/2\le\tfrac12$), there is $\tilde f$ with
\[
  a=E\tilde f,\qquad\norm{\pcw\tilde f}\le C\,\pcNorm a .
\]

\medskip\noindent\emph{Step 2 (three maps, from the data to the vectors of the proof).} Each of the following maps is linear on $L^2(\Hb)$, and is bounded
there and at the weight $\pcXzero$ on functions supported in one shell. Lemma~\ref{os:transfer} then bounds it at the weight $\pcw$.

\smallskip\noindent(a) \emph{Words.} For a word $Z_J$ of length $n$ let $T_ah:=Z_JB'_{\mathbf s}Eh$. Expand $B'_{\mathbf s}$ as in the proof of
Lemma~\ref{lem:wb}. Each term is a coefficient times a word of length at most $L$ applied to $Eh$, and by Lemma~\ref{os:words}(b), with $\gamma=1$ and
$m=3$ (so $2(\gamma+\tfrac12)^2<m^2$), it is bounded at the weight $\pcXzero$ by $2(L+1)L!\,(QC_0/\sigma)^L\norm{\pcXzero h}$. This is the bound of Lemma~\ref{lem:sc}
with $(2,QC_0,\sigma)$ in place of $(B_\beta,C_3,\alpha)$, so the count in the proof of Lemma~\ref{lem:wb} gives, for a suitable $C_5\ge C_4$,
\[
  \norm{\pcXzero Z_JB'_{\mathbf s}Eu}\le F_n\norm{\pcXzero u},\qquad\norm{Z_JB'_{\mathbf s}Eu}\le F_n\norm u ;
\]
the second bound is the same argument without a weight. Lemma~\ref{os:transfer} with $M_0=M_1=F_n$ gives
\begin{equation}\label{pc:wbw}
  \norm{\pcw Z_JB'_{\mathbf s}Eh}\le K\,F_n\,\norm{\pcw h}.
\end{equation}
In particular every word applied to $a=E\tilde f$ lies in $L^2_{\pcw}$.

\smallskip\noindent(b) \emph{Flow.} For $s',s''\ge0$ with $\hbar|s'-s''|\le\varepsilon$, the map $V_\varepsilon(s',s'')$ has norm $1$ on $L^2$ and, by
Proposition~\ref{prop:ws}, norm at most $\ue^{G_1/2}$ at the weight $\pcXzero$. So
\begin{equation}\label{pc:flow}
  \norm{\pcw V_\varepsilon(s',s'')h}\le K\ue^{G_1/2}\norm{\pcw h}.
\end{equation}

\smallskip\noindent(c) \emph{Defect.} For $0\le t'\le s_1\le t$ let
\[
  T_ch:=\bigl(\Delta(t')-\Delta_\varepsilon(t')\bigr)V_\varepsilon(t',s_1)B'_{\mathbf s}Eh .
\]
\emph{Unweighted bound.} By Lemma~\ref{os:words}, $\norm{\ue^{\sigma\Lambda}Eh}\le\norm h$. By (WB-u) of Lemma~\ref{lem:wb}, with $\sigma$ in the role of
$\alpha$ and the radii $a_1:=\tfrac34\sigma<\sigma$, by Proposition~\ref{prop:a0} from $a_1$ to $a_2:=\tfrac12\sigma$ (valid for
$t\le\tfrac13c_A\varepsilon/\hbar$), and by Lemma~\ref{lem:unw} with $a=a_2$, where $a_2/(8\varepsilon)=x/8$,
\[
  \norm{T_ch}\le M_0\norm h,\qquad M_0:=C\sigma^{-4}\ue^{-x/8}\,C_4^{N+K+1}(2N)!\,(2K+2)!\,\sigma^{-(2K+2)}(\sigma\hbar)^{-2N}.
\]
\emph{Weighted bound.} Let $u$ be supported in one shell and assume $\hbar s_1\le\varepsilon$. By (a), $\pcXzero ZB'_{\mathbf s}Eu\in L^2$ for $Z\in\{1,X,Y,X^2,Y^2\}$, and
Lemma~\ref{lem:path}(b) with $\gamma=1$ and $b'=B'_{\mathbf s}Eu$ gives
\[
  \norm{\pcXzero T_cu}\le M_1\norm{\pcXzero u},\qquad M_1:=C\varepsilon^{-3}F_2 .
\]
Lemma~\ref{os:transfer} requires $M_0\le M_1$; we apply it with $M_0+M_1$ in place of $M_1$. Since $M_0\le C\sigma^{-4}\ue^{-x/8}F_2$,
\begin{equation}\label{pc:defect}
  \norm{\pcw T_ch}\le K\bigl(M_0+M_0^{1/4}M_1^{3/4}\bigr)\norm{\pcw h}\le C\,\varepsilon^{-3}\sigma^{-4}\,\ue^{-x/32}\,F_2\,\norm{\pcw h}.
\end{equation}

\medskip\noindent\emph{Step 3 (Duhamel on $M$ for data of bounded support).} Let $h\in L^2$ be supported in $\{\psi\le R_0\}$ for some $R_0$, and let
$b'_h:=B'_{\mathbf s}Eh$, $g(t'):=V_\varepsilon(t',s_1)b'_h$ for $0\le t'\le s_1$, $w_h:=g(0)$ and $d_h(t'):=T_ch$. Then $\pcXzero h\in L^2$ and, by (a), all
words applied to $b'_h$ lie in $L^2(\pcXzero^2d\nu)$. The weight $\pcXzero^{3/4}$ satisfies (W1) and (W2) (Lemma~\ref{lem:t2}), and by
Lemma~\ref{lem:path}(c) with $\gamma=\tfrac34$ the path $g$ satisfies the hypotheses of Lemma~\ref{lem:surf-duh} for it. As in the proof of
Theorem~\ref{thm:dd-appx},
\[
  \mathcal V_{s_1}b'_h=\mathcal V_0w_h-\tfrac{i\hbar}2\int_0^{s_1}\mathcal V_{t'}d_h(t')\,dt' .
\]
By Remark~\ref{os:indep} each term may be estimated through the weight $\pcw$. By Proposition~\ref{os:p1}, \eqref{pc:flow} and \eqref{pc:wbw} with $n=0$,
\[
  \norm{\mathcal V_0w_h}\le C\norm{\pcw V_\varepsilon(0,s_1)b'_h}\le C\norm{\pcw b'_h}\le C\,F_0\,\norm{\pcw h},
\]
and by Proposition~\ref{os:p1} and \eqref{pc:defect},
\[
  \norm{\mathcal V_{t'}d_h(t')}\le C\ue^{t'/2}(1+t'/\log2)^{c'}\norm{\pcw d_h(t')}\le C\,\ue^{t'}\varepsilon^{-3}\sigma^{-4}\ue^{-x/32}F_2\norm{\pcw h}.
\]
Hence $\norm{\mathcal V_{s_1}b'_h}\le C(F_0+R\,F_2)\norm{\pcw h}$ with
\[
  R:=\hbar t\,\ue^{t}\,\varepsilon^{-3}\sigma^{-4}\,\ue^{-x/32}.
\]
\emph{Absorption.} Let $t\le\delta'x$ with $\delta'\le\tfrac1{64}$, so $\ue^t\le\ue^{x/64}$. Since $\sigma^3\ge\kappa'\hbar$ and $\sigma\le1$,
\[
  \sigma^{-1}\le x/\sqrt{\kappa'},\qquad\varepsilon^{-1}=2x/\sigma\le2x^2/\sqrt{\kappa'},\qquad\hbar t\le\delta'\sigma/x\le\delta'/x .
\]
So $R\le C(\kappa)\,x^{9}\ue^{-x/64}\le C(\kappa)\,x^{-2}$, and $x^{-2}=\hbar/\sigma\le\sigma^2/\kappa'$. Since $F_0\le\sigma^2F_2$,
\begin{equation}\label{pc:star}
  \norm{\mathcal V_{s_1}B'_{\mathbf s}Eh}\le C\,\sigma^2F_2\,\norm{\pcw h}.
\end{equation}

\medskip\noindent\emph{Step 4 (density).} The map $h\mapsto\mathcal V_{s_1}B'_{\mathbf s}Eh$ is linear and bounded from $L^2(\pcw^2d\nu)$ to $L^2(M)$, by
Proposition~\ref{os:p1} and \eqref{pc:wbw} with $n=0$. The functions of bounded $\psi$-support are dense in $L^2(\pcw^2d\nu)$. So \eqref{pc:star} holds for
every $h\in L^2(\pcw^2d\nu)$, in particular for $\tilde f$:
\[
  \norm{\mathcal V_{s_1}b'(\mathbf s)}\le C\,C_5^{N+K+3}(2N)!\,(2K+4)!\,\sigma^{-(2K+2)}(\sigma\hbar)^{-2N}\norm{\pcw\tilde f}.
\]

\medskip\noindent\emph{Step 5 (collecting).} By \eqref{pc:wbw}, every word applied to $a$ lies in $L^2_{\pcw}$, and $\pcw$ satisfies (W1), (W2). So (D4) of
Proposition~\ref{prop:dd-id} holds with the weight $\pcw$: the error is at most $(\hbar/2)^{K+1}$ times the integral of $\norm{\mathcal V_{s_1}b'(\mathbf s)}$ over
$\Sigma_{K+1}(t)$. As in the proof of Theorem~\ref{thm:dd-appx}, with $(2K+4)!/(K+1)!\le C^{K+1}K!$, $(2N)!\le4^NN^{2N}$ and
$\hbar^{K+1}t^{K+1}\sigma^{-(2K+2)}=(\hbar t/\sigma^2)^{K+1}$, Step~4 and Step~1 give the bound of the theorem with $\sigma$ in place of $\alpha$. Since
$\sigma=\alpha/6$, $(\sigma\hbar)^{-2N}(\hbar t/\sigma^2)^{K+1}=6^{2N+2K+2}(\alpha\hbar)^{-2N}(\hbar t/\alpha^2)^{K+1}$, and the factor is absorbed in $C^{N+K+1}$.

\medskip\noindent\emph{Time windows.} Steps~2(b), 2(c) and~3 use $\hbar s_1\le\varepsilon$, that is $t\le x/2$; Proposition~\ref{prop:a0} uses
$t\le\tfrac13c_A\varepsilon/\hbar=c_Ax/6$; the absorption uses $t\le x/64$. So the proof holds for $t\le\delta'x$ with $\delta':=\min(\tfrac1{64},c_A/6)$, and
$\delta'x=(\delta'/\sqrt6)(\alpha/\hbar)^{1/2}$. The hypothesis $\alpha\le\tfrac12$ ensures $\sigma<\alpha\le1$ in Proposition~\ref{os:data} and keeps all
radii below $\tau_0=1$. \qed

\subsection{Corollaries}\label{subsec:pc-cor}
\begin{corollary}[Corollary 1.5 for the printed class]\label{pc:cor15}
Under the hypotheses of Theorem~\ref{pc:main}, put $\delta_K:=C\hbar t/\alpha^2$ and let $\epsilon\in(0,1)$, with $\delta\le\kappa^{1/2}/\max(1,C)$.
For every integer $K$ with $(1-\epsilon)/\delta_K\le K\le(1-\epsilon/2)/\delta_K$,
\[
  \norm{\Delta^N\big[\cU(t)\big(a\,\ue^{i\varphi_b/\hbar}\big)_\Gamma-u^{(K)}(t)_\Gamma\big]}_{L^2(M)}
  \le C'\epsilon^{-1/2}C^{N+1}N^{2N}(\alpha\hbar)^{-2N}\ue^{-(1-\epsilon)\alpha^2/(C\hbar t)}\pcNorm{a},
\]
with $C'$ absolute. If the window contains no integer the statement is vacuous; at $t=0$ the error is $0$.
\end{corollary}
\begin{proof}
The proof of Corollary~\ref{cor:dd15-appx} uses only the form of the bound of Theorem~\ref{thm:dd-appx}:
$C^{N+K+1}N^{2N}K!(\alpha\hbar)^{-2N}(\hbar t/\alpha^2)^{K+1}$ times the norm of $a$, for every $K$. Theorem~\ref{pc:main} gives a bound of
exactly this form with $\pcNorm{a}$ in place of the swapped norm, so that proof applies word for word.
\end{proof}

\begin{lemma}[local Sobolev inequality on a Fuchsian quotient]\label{pc:g43}
Let $m(z):=\#\{\gamma\in\Gamma:\ d(\tilde z,\gamma\tilde z)<2\}$ for $z\in M$ with a lift $\tilde z$; $m$ is finite and does not depend on the lift.
There is an absolute constant $C$ such that every $u\in\operatorname{Dom}(\Delta_M)$ has a continuous representative, and for all $z\in M$,
\[
  |u(z)|\le C\,m(z)^{1/2}\big(\norm{\Delta u}_{L^2(M)}+\norm{u}_{L^2(M)}\big).
\]
\end{lemma}
\begin{proof}
\emph{Step 1 (the lift).} Let $\tilde u$ and $\tilde f$ be the $\Gamma$-invariant lifts of $u$ and of $f:=\Delta_Mu$. Then $\Delta\tilde u=\tilde f$ in the
sense of distributions on $\Hb$: for $\phi\in C_c^\infty(\Hb)$, unfolding gives
$\langle\tilde u,\Delta\phi\rangle_{\Hb}=\langle u,(\Delta\phi)_\Gamma\rangle_M=\langle u,\Delta_M\phi_\Gamma\rangle_M=\langle f,\phi_\Gamma\rangle_M=\langle\tilde f,\phi\rangle_{\Hb}$.

\emph{Step 2 (counting).} For a $\Gamma$-invariant $\tilde h$ and a fundamental domain $F$,
$\int_{B(\tilde z,1)}|\tilde h|^2d\nu=\int_Fn(w)|h(w)|^2d\nu(w)$, where $n(w):=\#\{\gamma:\gamma w\in B(\tilde z,1)\}$. If $\gamma_0w$ and $\gamma w$ both lie in
$B(\tilde z,1)$, then $d(\tilde z,\gamma\gamma_0^{-1}\tilde z)<2$; hence $n(w)\le m(z)$ and $\int_{B(\tilde z,1)}|\tilde h|^2\le m(z)\norm{h}^2_{L^2(M)}$. In
particular $\tilde u$ and $\tilde f$ are locally square integrable. The number $m(z)$ is finite because $\Gamma$ acts properly discontinuously.

\emph{Step 3 (interior regularity).} Let $v\in L^2(B(\tilde z,1))$ with $\Delta v\in L^2(B(\tilde z,1))$ in the sense of distributions. By interior
elliptic regularity, uniform in $\tilde z$ by the homogeneity of $\Hb$, $v\in H^2(B(\tilde z,\tfrac12))$ with
$\norm v_{H^2(B(\tilde z,1/2))}\le C(\norm{\Delta v}_{L^2(B(\tilde z,1))}+\norm v_{L^2(B(\tilde z,1))})$, and since $H^2\subset C^0$ in dimension two, $v$ is
continuous on $B(\tilde z,\tfrac12)$ with $|v(\tilde z)|\le C\norm v_{H^2(B(\tilde z,1/2))}$.

Apply Step~3 to $v=\tilde u$ and use Step~2 for $\tilde h=\tilde u$ and $\tilde h=\tilde f$. The continuous representative of $\tilde u$ is $\Gamma$-invariant
and descends to $M$. \qed
\end{proof}

\begin{corollary}[Corollary 1.6 on bounded geometry, printed class]\label{pc:cor16bg}
Under the hypotheses of Theorem~\ref{pc:main}, assume in addition that the injectivity radius of $M$ is at least $r_0>0$. Then for every $z\in M$,
\[
  \big|\Delta^N\big[\cU(t)\big(a\,\ue^{i\varphi_b/\hbar}\big)_\Gamma-u^{(K)}(t)_\Gamma\big](z)\big|
  \le C^{N+K+1}(2N)!\,K!\,(\alpha\hbar)^{-2N-4}\Big(\frac{\hbar t}{\alpha^2}\Big)^{K+1}\pcNorm{a},
\]
with $C$ depending also on $r_0$.
\end{corollary}

\begin{corollary}[Corollary 1.6 on every Fuchsian quotient, with the cusp factor]\label{pc:cor16}
Under the hypotheses of Theorem~\ref{pc:main}, for every $z\in M$,
\begin{multline*}
  \big|\Delta^N\big[\cU(t)\big(a\,\ue^{i\varphi_b/\hbar}\big)_\Gamma-u^{(K)}(t)_\Gamma\big](z)\big|\\
  \le C\,m(z)^{1/2}\,C^{N+K+1}(2N)!\,K!\,(\alpha\hbar)^{-2N-4}\Big(\frac{\hbar t}{\alpha^2}\Big)^{K+1}\pcNorm{a}.
\end{multline*}
In both corollaries the error lies in $\operatorname{Dom}(\Delta_M^N)$, and the left side is the value of the continuous representative.
\end{corollary}
\begin{proof}
Let $E$ be the bracket in the statement. By Theorem~\ref{pc:main} at $N$ and at $N+1$, $u:=\Delta^NE$ lies in $\operatorname{Dom}(\Delta_M)$. Apply
Lemma~\ref{pc:g43} to $u$, and use $N^{2N}\le(\ue/2)^{2N}(2N)!$, $(N+1)^{2N+2}\le54^{N+1}(2N)!$ and $\alpha\hbar\le1$. This proves
Corollary~\ref{pc:cor16}. Under the hypothesis of Corollary~\ref{pc:cor16bg}, $m(z)\le C(r_0)$, as in the proof of Corollary~\ref{cor:dd16-appx}. \qed
\end{proof}

\begin{remark}[cusps]\label{pc:remcusp}
On a bounded-geometry quotient $m(z)\le C(r_0)$, so Corollary~\ref{pc:cor16bg} is the case of Corollary~\ref{pc:cor16}. In an embedded cusp of
height $H$, deep enough, $m(z)\asymp H$, so the factor $m(z)^{1/2}$ grows like the square root of the cusp height. Lemma~\ref{pc:g43} is sharp in this sense:
the cusp bump $u_S$ of Remark~\ref{rem:g4} at height $H=\ue^{S}$ has $\sup|u|/(\norm{\Delta u}+\norm{u})\asymp H^{1/2}$.

The proofs of the two corollaries use one Laplacian and give the exponent $-2N-2$ of $\alpha\hbar$; we state them with $-2N-4$ for the comparison with
\cite[(1.34)]{S}. Since $m(z)^{1/2}(\alpha\hbar)^{-2N-4}\le(\alpha\hbar)^{-2N-8}$ when $m(z)\le(\alpha\hbar)^{-8}$, Corollary~\ref{pc:cor16} gives the printed
prefactor $(\alpha\hbar)^{-2N-8}$ of \cite[(1.34)]{S} at every such $z$, with $(\hbar t/\alpha^2)^{K+1}$ in place of $(\hbar t)^K$ as in
Corollary~\ref{pc:cor16bg} (the literal $(\hbar t)^K$ only for fixed $\alpha$). In a cusp this is the part below height $\asymp(\alpha\hbar)^{-8}$.
Above that height the factor cannot be removed in general: on $\mathrm{PSL}(2,\mathbb Z)\backslash\Hb$, for $b$ a cusp, Theorem~\ref{cl:main} gives a lower
bound of the same order $m(z)^{1/2}$ up to powers of $\hbar$. The proof of Lemma~\ref{pc:g43} also gives the local
form $|u(z)|\le C\,m(z)^{1/2}\bigl(\norm{\Delta u}_{L^2(B_M(z,1))}+\norm{u}_{L^2(B_M(z,1))}\bigr)$, with $B_M(z,1)$ the image of $B(\tilde z,1)$. So the
factor-free bound would follow if, for $u=\Delta^NE$, both $\norm{u}_{L^2(B_M(z,1))}$ and $\norm{\Delta u}_{L^2(B_M(z,1))}$ were at most
$C\,m(z)^{-1/2}(\alpha\hbar)^{-p}$ times the global bounds at $N$ and $N+1$ (with $p\le6$ for the printed prefactor). Theorem~\ref{cl:main} shows that this
condition fails high in the cusp of the modular surface.
\end{remark}

% Subsection for sec_printed.tex: the cusp factor of Corollary pc:cor16 is necessary.
% Notes source: notes/Cor16_cusps.md, Propositions G4.4 and G4.4P (reviews R22, R22c, R23, R23b). Labels prefix cl.
\subsection{The cusp factor is necessary}\label{subsec:cl}

Throughout this subsection $\Gamma=\mathrm{PSL}(2,\mathbb Z)$, $M=\Gamma\backslash\Hb$ (cusp at $\infty$ of width $1$), $x_0=i$, $X_0=\cosh d(i,\cdot)$,
$b=\infty$ and $\varphi=-\log y$. For a function $F$ on $\Hb$ we write $[F]_\Gamma(z)=\sum_{\gamma\in\Gamma}F(\gamma z)$, and
$E_K(t):=\cU(t)(a\,\ue^{i\varphi/\hbar})_\Gamma-u^{(K)}(t)_\Gamma$. Words $Z_S$ are products of $X$ and $Y$.

\begin{theorem}[the cusp factor is necessary]\label{cl:main}
Let $\beta>\tfrac12$, let $\alpha_0<\tfrac12$ be small enough for Theorems~\ref{thm:dd-appx} and~\ref{pc:main} and for Lemma~\ref{cl:amp}(i), and fix
$\alpha\in(0,\alpha_0]$; the constants below depend on the fixed $\alpha$, $\beta$ and $\kappa$.
\begin{enumerate}
\item[(i)] (Swapped class.) For $a=\ue^{-20X_0}$ there are $\hbar_0,c_*,t_*>0$ such that for $0<\hbar\le\hbar_0$, every
$t$ with $\max(2\log(1/\hbar),t_*)\le t\le\min(\hbar^{-1/5},\delta(\alpha/\hbar)^{1/2})$, and every $z=x+i\ue^t$,
\[
|E_0(t)(z)|\ge c_*\,\ue^{t/2}\hbar^{3/2}t,\qquad m(z)\le5\ue^t .
\]
\item[(ii)] (Printed class.) The same holds for $a=\ue^{-\alpha\Lambda}\ue^{-\kappa X_0}$, with $\kappa\ge\kappa_*$ and $\alpha\le\alpha_*(\kappa):=\min(\alpha_0,c\kappa^{-6})$; here the
proof uses Theorem~\ref{pc:main}.
\end{enumerate}
Consequently, on $\mathrm{PSL}(2,\mathbb Z)\backslash\Hb$ and for this $b$, the factor $m(z)^{1/2}$ in Corollary~\ref{pc:cor16} (and in its swapped-class
version) cannot be removed. For every $p\ge0$ no bound $\sup_M|E_0(t)|\le C(\alpha\hbar)^{-p}\norm a$ with $C$ independent of $\hbar$ holds on the window (take
$t=(2p+4)\log(1/\hbar)$), and for $\gamma<\tfrac12$ no bound $C\,m(z)^{\gamma}\hbar^{-p}\norm a$ with $C$ independent of $\hbar$ holds (take $t=\hbar^{-1/5}$). In particular the printed Corollary~1.6, \cite[(1.34)]{S} at
$N=K=0$, which is claimed for every Fuchsian $\Gamma$ and every $b$, is false when its constants $C,\delta$ do not depend on $\hbar$
(Remark~\ref{cl:reading}): at $t=20\log(1/\hbar)$ the lower bound (ii) exceeds $C(\alpha\hbar)^{-8}\norm a_{\alpha,\beta}$ for small $\hbar$.
\end{theorem}

\begin{remark}[the reading of the printed constants]\label{cl:reading}
\cite[Theorem~1.4 and Corollary~1.6]{S} choose $C,\delta$ after ``for all $\alpha,\beta$ with $\beta>1/2$ and $\alpha^3\ge\beta^2\hbar$''
and do not say whether they may depend on $\hbar$. We read them as independent of $\hbar$, as in a semiclassical estimate and as in all
our theorems. Under that reading the printed Corollary~1.6 is false. The contradiction is asymptotic: it needs one $C$ valid as $\hbar$ tends
to $0$. It does not rule out a bound with $\hbar$-dependent constants, and we do not claim that such a bound holds.
\end{remark}

\begin{remark}[the approximation is still accurate there]\label{cl:relative}
At the points of Theorem~\ref{cl:main}, in case (i) (the swapped-class amplitude $a=\ue^{-20X_0}$), the solution itself is large and the error remains
small relative to it: for
$\hbar\le\hbar_0$ and $t$ in the range of Theorem~\ref{cl:main},
\[
c\,\hbar t\le\frac{|E_0(t)(z)|}{|u^{(0)}(t)_\Gamma(z)|}\le C\,\hbar t,\qquad z=x+i\ue^t .
\]
So what fails is the form of the printed bound (a fixed multiple of $\norm a$, independent of $z$), not the accuracy of the approximation: the
periodized approximation and the solution both grow like $m(z)^{1/2}\hbar^{1/2}$ at these points, and their difference is a factor $O(\hbar t)$
smaller. (This is a statement about this amplitude on the window of Theorem~\ref{cl:main}; it is not a general relative-error theorem.)
\end{remark}
\begin{proof}
Lemma~\ref{cl:asym} with $A=a$ gives
\[
u^{(0)}(t)_\Gamma(z)=\ue^{-it/(2\hbar)}\bigl(\ue^{t/2}\sqrt{\pi\hbar}\,\ue^{i\pi/4}\ue^{it/\hbar}S_a(\hbar)+O(\ue^{t/2}\hbar^{3/2})\bigr),
\]
where $S_a(\hbar)=\sum_{c\ge1}c^{2i/\hbar}\sum_{(a',c)=1}a(a'/c+i/c^2)$.
With $a=\ue^{-\kappa X_0}$ and $\kappa=20$, $a(i)=\ue^{-\kappa}$ and the other terms are at most $\ue^{-\kappa}\ue^{-\kappa(X_0-1)}$; the count in the proof of
Lemma~\ref{cl:amp}(iii) gives $0.99\,\ue^{-\kappa}\le|S_a(\hbar)|\le1.01\,\ue^{-\kappa}$. The proof of Theorem~\ref{cl:main} gives
\[
|E_0(z)|=\ue^{t/2}\sqrt{\pi\hbar}\,\tfrac{\hbar t}2\,|S_0(\hbar)|\,\bigl(1+O(1/t)+O(\hbar t)+O(\hbar^{1/2}t^2)\bigr),
\]
and $0.98\,\kappa\ue^{-\kappa}\le|S_0(\hbar)|\le1.02\,\kappa\ue^{-\kappa}$. Divide.
\end{proof}

\begin{remark}[scope]\label{cl:scope}
Theorem~\ref{cl:main} extends to every $b\in\mathbb Q$, the other cusps of $\mathrm{PSL}(2,\mathbb Z)$: if $b=\gamma\infty$ with $\gamma\in\Gamma$, then
$\varphi_b=\varphi_\infty\circ\gamma^{-1}$ up to a constant, the approximation for $b$ is the $\gamma$-transport of the one for $\infty$, and the error for
$(b,a)$ is a constant phase times the error for $(\infty,a\circ\gamma)$. So take $a_b:=a\circ\gamma^{-1}$, that is $\ue^{-\kappa\cosh d(\gamma i,\cdot)}$ in case
(i) and $\ue^{-\alpha\Lambda}\ue^{-\kappa\cosh d(\gamma i,\cdot)}$ in case (ii) ($\Lambda$ commutes with isometries); these have finite norms about the base
point $i$, since $\langle d(i,\cdot)\rangle\le\langle d(\gamma i,\cdot)\rangle+d(i,\gamma i)$ and the tube argument of Lemma~\ref{cl:amp}(i) is isometry
invariant. Like the printed Corollary~1.6, Corollary~\ref{pc:cor16} is stated uniformly for every Fuchsian $\Gamma$ and every $b$; by Theorem~\ref{cl:main} no
statement of that form can have $m(z)^\gamma$ with $\gamma<\tfrac12$, so the exponent $\tfrac12$ is optimal among them (the power of $\hbar$ is not
claimed optimal). For irrational $b$, and for cusped quotients other than $\mathrm{PSL}(2,\mathbb Z)\backslash\Hb$, the necessity of the factor is not decided
here; a bound depending on $b$ without that factor is not excluded. The mechanism below is not special to $b=\infty$: in the frame where
$b=\infty$, the top of every $\Gamma$-image of a cusp horocycle is a stationary point of $-\log y$.
\end{remark}

The mechanism is the following. The error contains the explicit first Dyson term $\ue^{-it/(2\hbar)}(S(t)\,\tfrac{i\hbar}2P_1(t)a)\ue^{i\varphi/\hbar}$,
periodized. High in the cusp the sum over $\Gamma$ concentrates, by stationary phase, at the tops $a/c+i/(c^2Y)$ of the horocycles $\gamma(\mathbb R+iY)$,
whose images under $D_t:(x,y)\mapsto(x,\ue^ty)$ are the points $a/c+i/c^2$,
and its size is $\ue^{t/2}\sqrt{\pi\hbar}\cdot\hbar t$; the remainder is controlled by the $L^2$ theorem.

\begin{lemma}[semiclassical local Sobolev inequality]\label{cl:ssob}
For every Fuchsian $\Gamma$ there is an absolute $C$ with $|u(z)|\le C\,m(z)^{1/2}\bigl(r^{-1}\norm u_{L^2(M)}+r^3\norm{\Delta^2u}_{L^2(M)}\bigr)$ for
$0<r\le1$, $z\in M$ and $u\in H^4(M)$. For $\Gamma=\mathrm{PSL}(2,\mathbb Z)$ and $z=x+iY$ with $Y\ge4$, $m(z)\le5Y$.
\end{lemma}
\begin{proof}
By homogeneity take $\tilde z=i$. In normal coordinates $\xi$ at $i$ put $v(\xi)=\tilde u(\exp_i(r\xi))$ for $|\xi|<1$. Then
$r^4(\Delta^2\tilde u)(\exp_i(r\xi))=L_r^2v(\xi)$, where $L_r$ is $r^2\Delta$ in these coordinates: a uniformly elliptic operator whose coefficients are
bounded in $C^\infty$ uniformly for $0<r\le1$. Interior regularity for $L_r^2$ and $H^4\subset L^\infty$ give
$|v(0)|\le C(\norm v_{L^2(B_1)}+\norm{L_r^2v}_{L^2(B_1)})$. Since $d\nu\asymp r^2d\xi$ on the ball,
$|\tilde u(i)|\le C(r^{-1}\norm{\tilde u}_{L^2(B(i,r))}+r^3\norm{\Delta^2\tilde u}_{L^2(B(i,r))})$, and $B(\tilde z,r)\subset B(\tilde z,1)$, so the unfolding step of
Lemma~\ref{pc:g43} applies. For the count: if $\gamma\notin\Gamma_\infty$ then $\operatorname{Im}\gamma z\le1/Y$ and $d(z,\gamma z)\ge2\log Y\ge2$; and
$d(z,z+n)=2\operatorname{arcsinh}(|n|/(2Y))<2$ if and only if $|n|<2Y\sinh1$. So $m(z)\le4Y\sinh1+1\le5Y$ for $Y\ge4$.
\end{proof}

\begin{lemma}[cusp asymptotics]\label{cl:asym}
Let $\vartheta$ be either $X_0^{-2}$ or $\ue^{-d(i,\cdot)}(1+d(i,\cdot))^{-3/2}$, and for $A\in C^4(\Hb)$ put
$B(A):=\sup_w\vartheta(w)^{-1}\sum_{|S|\le4}|Z_SA(w)|$. There is a $C$ (depending only on the choice of $\vartheta$) such that for $0<\hbar\le\tfrac12$,
$t\ge2\log(1/\hbar)$, $Y=\ue^t$ and $z=x+iY$,
\begin{multline*}
\Bigl|\bigl[(S(t)A)\ue^{i\varphi/\hbar}\bigr]_\Gamma(z)-\ue^{t/2}\sqrt{\pi\hbar}\,\ue^{i\pi/4}\ue^{it/\hbar}\sum_{c\ge1}c^{2i/\hbar}\sum_{a\in\mathbb Z,(a,c)=1}A\Bigl(\frac ac+\frac i{c^2}\Bigr)\Bigr|\\
\le C\,\ue^{t/2}\hbar^{3/2}B(A),
\end{multline*}
and the double series converges absolutely.
\end{lemma}
\begin{proof}
\emph{Cosets.} An element of $\mathrm{PSL}(2,\mathbb Z)$ is a translation $T^k$, or has lower-left entry $c\ge1$ and acts by $w\mapsto a/c-1/(c(cw+d))$ with
$ad\equiv1\pmod c$. For fixed $c$ and $a$ coprime to $c$, $d$ runs over $d_a+c\mathbb Z$, and $d=d_a+cn$ corresponds to $w=z+n$. Hence
$[F]_\Gamma(z)=\sum_kF(z+k)+\sum_{c\ge1}\sum_{(a,c)=1}\sum_{n\in\mathbb Z}G_{a,c}(n)$ with $G_{a,c}(u)=F(\gamma_{a,c}(z+u))$. The translation part is
bounded by $\ue^{-t/2}B(A)\sum_k\vartheta(x+k+i\ue^{2t})\le C\ue^{-t/2}B(A)$, which is $\le C\ue^{t/2}\hbar^{3/2}B(A)$.

\emph{Change of variable.} With $\tau=(c(x+u)+d_a)/Y$, the point $D_t\gamma_{a,c}(z+u)$, where $D_t(x,y)=(x,\ue^ty)$ and $S(t)A=\ue^{-t/2}A\circ D_t$, is
\[
Q(\tau)=\frac ac-\frac{\tau}{cY(\tau^2+c^2)}+\frac{i}{\tau^2+c^2},\qquad \varphi(\gamma_{a,c}(z+u))=t+\log(\tau^2+c^2),
\]
so $G_{a,c}(u)=g(\tau):=\ue^{-t/2}A(Q(\tau))\,\ue^{i(t+\log(\tau^2+c^2))/\hbar}$ and $du=(Y/c)\,d\tau$. Along $Q$,
$\partial_\tau A(Q)=-(YA)(Q)\,\frac{2\tau}{\tau^2+c^2}-(XA)(Q)\,\frac{c^2-\tau^2}{cY(\tau^2+c^2)}$, so $|\partial_\tau^jA(Q)|\le C_jc^{-j}B(A)\vartheta(Q)$ for
$j\le4$. From $X_0(Q)=(q^2+1)/(2\operatorname{Im}Q)+\operatorname{Im}Q/2$ with $|q-a/c|\le1/(2c^2Y)$ one gets
$\vartheta(Q(\tau))\le C\,\vartheta_{a,c}\,(1+\tau^2/c^2)^{-s}$, where $s=2$, $\vartheta_{a,c}=(a^2+c^2)^{-2}$ for $\vartheta=X_0^{-2}$, and $s=1$,
$\vartheta_{a,c}=(a^2+c^2)^{-1}(1+\log(a^2+c^2))^{-3/2}$ for the second choice. In both cases $\sum_{(a,c)}\vartheta_{a,c}<\infty$ (for the second, by
partial summation, since $\#\{(a,c):a^2+c^2\le N\}\le CN$).

\emph{Poisson summation.} $\sum_nG(n)=\sum_m\hat G(m)$ with
\[
\hat G(m)=\frac Yc\,\ue^{2\pi im(x+d_a/c)}\int g(\tau)\,\ue^{-2\pi imY\tau/c}\,d\tau .
\]
For $m\ne0$ the total phase $\Phi$ has $|\Phi'|\ge2\pi|m|Y/c-1/(\hbar c)\ge\pi|m|Y/c$ (as $\hbar Y\ge\hbar^{-1}\ge2$), $|\Phi''|\le2/(\hbar c^2)$ and
$|\Phi'''|\le C/(\hbar c^3)$. Two integrations by parts give
\[
\Bigl|\int g\,\ue^{-2\pi imY\tau/c}d\tau\Bigr|\le C\Bigl(\frac c{|m|Y}\Bigr)^2c^{-2}\ue^{-t/2}B(A)\int\vartheta(Q(\tau))\,d\tau\le C\ue^{-t/2}B(A)\,\vartheta_{a,c}\,\frac c{m^2Y^2},
\]
using $\int(1+\tau^2/c^2)^{-s}d\tau\le Cc$. Hence $\sum_{m\ne0}|\hat G(m)|\le C\ue^{-t/2}B(A)\vartheta_{a,c}/Y$, and the sum over $(a,c)$ is $\le C\ue^{-3t/2}B(A)$.

\emph{Stationary phase.} With $\tau=c\sigma$, $\hat G(0)=Y\ue^{it/\hbar}c^{2i/\hbar}\int f(\sigma)\ue^{i\theta(\sigma)/\hbar}d\sigma$, where
$f(\sigma)=\ue^{-t/2}A(Q(c\sigma))$ and $\theta(\sigma)=\log(1+\sigma^2)$. Let $\chi$ be a cutoff, $\chi=1$ on $[-\tfrac12,\tfrac12]$, $\operatorname{supp}\chi\subset[-1,1]$.
Near part: $\theta=s(\sigma)^2$ with $s(\sigma)=\operatorname{sign}(\sigma)(\log(1+\sigma^2))^{1/2}$ smooth and $s'(0)=1$; for $h(s)=(\chi f)(\sigma(s))\sigma'(s)$,
\[
\int h\,\ue^{is^2/\hbar}ds=\sqrt{\pi\hbar}\,\ue^{i\pi/4}\,\frac1{2\pi}\int\hat h(\xi)\,\ue^{-i\hbar\xi^2/4}d\xi,
\]
and $|\ue^{-i\hbar\xi^2/4}-1|\le\hbar\xi^2/4$ gives the error $C\hbar^{3/2}\norm h_{C^4}\le C\hbar^{3/2}\ue^{-t/2}B(A)\vartheta_{a,c}$. Far part, on $|\sigma|\ge\tfrac12$, where $|\theta'|\ge(1+|\sigma|)^{-1}/2$ and
$|\theta^{(j)}|\le C(1+|\sigma|)^{-j}$. First truncate: with a cutoff $\chi_Y$ equal to $1$ on $|\sigma|\le Y$ and $0$ on $|\sigma|\ge2Y$, the integral over
the support of $1-\chi_Y$ is bounded absolutely by $C\ue^{-t/2}B(A)\vartheta_{a,c}/Y\le C\hbar^2\ue^{-t/2}B(A)\vartheta_{a,c}$. On $\tfrac12\le|\sigma|\le2Y$
integrate by parts once with $D^Tu=-(u/\theta')'$; the derivatives of $\chi_Y$ add terms of order $\hbar/Y$. In $\partial_\sigma f$ the $Y$-terms carry
$2\sigma/(1+\sigma^2)$ and gain $(1+|\sigma|)^{-1}$ per derivative, and the $X$-terms carry $1/Y$. The $X$-part of $D^T[(1-\chi)\chi_Yf]$ is integrated directly,
giving at most $C\hbar\,\ue^{-t/2}B(A)\vartheta_{a,c}Y^{-1}\log Y\le C\hbar^{3/2}\ue^{-t/2}B(A)\vartheta_{a,c}$ (as $Y^{-1}\log Y=t\ue^{-t}\le\hbar^{1/2}$); the
remaining symbol part is integrated by parts once more, giving at most $C\hbar^2\ue^{-t/2}B(A)\vartheta_{a,c}$. Since $f(0)=\ue^{-t/2}A(a/c+i/c^2)$,
\[
\hat G(0)=\ue^{t/2}\sqrt{\pi\hbar}\,\ue^{i\pi/4}\ue^{it/\hbar}c^{2i/\hbar}A\Bigl(\frac ac+\frac i{c^2}\Bigr)+O\bigl(\ue^{t/2}\hbar^{3/2}B(A)\vartheta_{a,c}\bigr).
\]
Summing over $(a,c)$ proves the lemma.
\end{proof}

\begin{lemma}[the Poisson kernel off the diagonal]\label{cl:kernel}
Fix $r_0\in(0,1]$. For $0<\alpha\le\alpha_0<\pi/2$ the kernel $k_\alpha(r)$ of $\ue^{-\alpha\Lambda}$ extends holomorphically to
$\{a+i\theta:a>3r_0/4,\ |\theta|<r_0/4\}$ with $|\tilde k_\alpha(a+i\theta)|\le C\alpha\,\ue^{-a}(1+a)^{-3/2}$. Hence
$|\partial_r^jk_\alpha(r)|\le C_j\alpha\,\ue^{-r}(1+r)^{-3/2}$ for $r\ge r_0$, uniformly in $\alpha\le\alpha_0$.
\end{lemma}
\begin{proof}
By Lemma~\ref{os:hk} with $\theta_0=r_0/4$,
\[
|\tilde k_\alpha(a+i\theta)|\le C\int_0^\infty\alpha(4\pi)^{-1/2}s^{-3/2}\ue^{-(\alpha^2-\theta^2)/(4s)}p_s(a)\,ds .
\]
The computation in the proof of Lemma~\ref{os:poisson}, with $\alpha^2-\theta^2$ in place of $\alpha^2$, turns this into
\[
C\alpha\int_a^\infty\frac{\rho\,K_2\bigl(\sqrt{\rho^2+\alpha^2-\theta^2}/2\bigr)}{\rho^2+\alpha^2-\theta^2}\,\frac{d\rho}{(\cosh\rho-\cosh a)^{1/2}},
\]
which converges because $\rho^2+\alpha^2-\theta^2\ge r_0^2/2$, and is bounded as in \textup{(\ref{os:ub})}. Cauchy estimates on discs of radius $r_0/8$ give the derivative bounds.
\end{proof}

\begin{lemma}[the amplitudes]\label{cl:amp}
(i) $a=\ue^{-\kappa X_0}$ has finite swapped norm for every $\alpha\le r_1<\tfrac12$ and $\beta$, and $|Z_Sa|\le C_{S,\kappa}X_0^{|S|}\ue^{-\kappa X_0}$.
(ii) $a=\ue^{-\alpha\Lambda}\ue^{-\kappa X_0}$ has $\norm a_{\alpha,\beta}=\norm{\ue^{\beta\langle d\rangle}\ue^{-\kappa X_0}}<\infty$, and
$|Z_Sa(w)|\le C_{S,\kappa}\ue^{-d(i,w)}(1+d(i,w))^{-3/2}$ uniformly in $\alpha\le\alpha_0$.
(iii) Put $A_0=(Y^2-Y)a$ and $S_0(\hbar)=\sum_{c\ge1}c^{2i/\hbar}\sum_{(a,c)=1}A_0(a/c+i/c^2)$. In case (i) with $\kappa=20$, $|S_0(\hbar)|\ge0.98\,\kappa\ue^{-\kappa}$
for all $\hbar$; in case (ii), $|S_0(\hbar)|\ge\tfrac12\kappa\ue^{-\kappa}$ for $\kappa\ge\kappa_*$, $\alpha\le c\kappa^{-6}$ and all $\hbar$.
\end{lemma}
\begin{proof}
(i) $f=\ue^{\beta\langle d\rangle}a$ continues to the crown with $\operatorname{Re}\tilde X_0=\cos(2r)X_0$ and $\operatorname{Re}\widetilde{\langle d\rangle}\le\langle d\rangle+C_Hr^2$ (proof of
Lemma~\ref{lem:ex}), so its tube norm is finite and $\norm{\ue^{\alpha\Lambda}f}\le C_0N_\alpha(f)$ as there. The vector fields act by $XX_0=x$, $YX_0=y-X_0$,
$Xx=y$, $Yy=y$, $Xy=Yx=0$, and $|x|,y\le2X_0$; this gives the word bound.
(ii) The norm is the definition. For the words, transport $i$ to $w'$ by the affine group, under which $X$ and $Y$ are invariant; then
$|Z_S^w\cosh d(w,w')|\le C_S\cosh d(w,w')$, and Lemma~\ref{cl:kernel} gives $|Z_S^wk_\alpha(d(w,w'))|\le C_S\alpha\ue^{-d(w,w')}(1+d(w,w'))^{-3/2}$ for
$d(w,w')\ge\tfrac14$. Split $h=\ue^{-\kappa X_0}$ near each $w$ as $h\psi_w+h(1-\psi_w)$, $\psi_w=1$ on $B(w,\tfrac18)$, $\operatorname{supp}\psi_w\subset B(w,\tfrac14)$:
the far part is bounded with the kernel estimate, and the near part (with $\psi_w$ frozen at the evaluation point) by local Sobolev and
$\norm{\ue^{-\alpha\Lambda}g}_{H^k}\le\norm g_{H^k}$; the words of $h$ carry powers of $\kappa$ and $X_0$, so the near part is at most
$C_{S,\kappa}\exp(-\tfrac\kappa2\cosh(\max(0,d(i,w)-\tfrac14)))$, which is dominated by $\ue^{-d}(1+d)^{-3/2}$.
(iii) $Y^2X_0=X_0$ gives $A_0=a[\kappa^2(y-X_0)^2-\kappa(2X_0-y)]$ in case (i), so $A_0(i)=-\kappa\ue^{-\kappa}$. At the other points $a/c+i/c^2$,
$X_0=(a^2+c^2+c^{-2})/2$, $X_0-1\ge\max(\tfrac12,(n-2)/2)$ with $n=a^2+c^2$, and $|A_0|\le\kappa\ue^{-\kappa}\ue^{-\kappa(X_0-1)}(\kappa X_0^2+2X_0)$. For
$\kappa=20$ and $q=\ue^{-10}<1/20000$: the two points $\pm1+i$ give at most $96q<0.005$; there is no coprime pair with $c\ge1$ and $n\in\{3,4\}$; the rest is
at most $\sum_{n\ge5}(4n+4)[5(n+1)^2+n+1]q^{n-2}\le5184q^3/(1-8q)<0.01$. In case (ii), $|A_0(i)+\kappa\ue^{-\kappa}|\le C\norm{(\ue^{-\alpha\Lambda}-1)h}_{H^4}\le
C\alpha\norm{(1-\Delta)^3h}\le C\alpha\kappa^6\ue^{-\kappa}$. The other points are $0.96$-separated ($\cosh d\ge\tfrac32$ for distinct ones). Split
$h=h_1+h_2$ with $h_1=h\psi$, $\psi=1$ on $B(i,0.44)$, $\operatorname{supp}\psi\subset B(i,0.45)$, so that $h_2$ vanishes on $B(i,0.44)$: $h_1$ contributes at most $C\alpha\kappa^{-1}\ue^{-\kappa}\vartheta(w)$ at each point by Lemma~\ref{cl:kernel}; for $h_2$, whose
derivatives up to order $4$ are at most $C(1+\kappa)^4(1+X_0)^4\ue^{-\kappa X_0}$, the split of (ii) at each point and the packing bound
$\#\{w:d(i,w)\le R\}=O(\ue^R)$ give at most $C\kappa^8\ue^{-\kappa\cosh0.44}$ in total, with $C$ independent of $\kappa\ge1$ and $\alpha\le\alpha_0$. Hence the sum
over the other points is at most $C(\alpha/\kappa^2+\kappa^8\ue^{-0.098\kappa})\kappa\ue^{-\kappa}$.
\end{proof}

\begin{proof}[Proof of Theorem~\ref{cl:main}]
With $c_\hbar=i\hbar/2$, the definition of $u^{(K)}$ gives $E_0=M_1+M_2+E_2$, where
$M_k=\ue^{-it/(2\hbar)}[(S(t)A_k)\ue^{i\varphi/\hbar}]_\Gamma$, $A_1=c_\hbar P_1(t)a$ and $A_2=c_\hbar^2P_2(t)a$. Since
$P_1(t)a=\int_0^t\Delta(s)a\,ds=\tfrac{1-\ue^{-2t}}2X^2a+t(Y^2-Y)a$, Lemma~\ref{cl:asym} gives
\[
|M_1(z)|\ge\ue^{t/2}\sqrt{\pi\hbar}\,\tfrac\hbar2\bigl(t|S_0(\hbar)|-S_X\bigr)-C\ue^{t/2}\hbar^{3/2}\hbar t,
\]
with $S_X=\sum_{(a,c)}|X^2a(a/c+i/c^2)|<\infty$. $P_2(t)a$ is a combination of
words of length at most $4$ with coefficients bounded by $t^2/2$, so $|M_2(z)|\le C\ue^{t/2}\hbar^{1/2}\hbar^2t^2$. For $E_2$, Theorem~\ref{thm:dd-appx} in case (i)
and Theorem~\ref{pc:main} in case (ii), at $K=2$ and $N=0,2$, give $\norm{E_2}\le C(\hbar t/\alpha^2)^3\norm a$ and
$\norm{\Delta^2E_2}\le C(\alpha\hbar)^{-4}(\hbar t/\alpha^2)^3\norm a$; Lemma~\ref{cl:ssob} with $r=\hbar$ gives $|E_2(z)|\le C\ue^{t/2}\hbar^{-1}\alpha^{-10}(\hbar t)^3\norm a$.
Hence
\[
|E_0(z)|\ge\ue^{t/2}\hbar^{3/2}t\Bigl[\tfrac{\sqrt\pi}2|S_0(\hbar)|-\tfrac{\sqrt\pi}2\frac{S_X}t-C\hbar-C\hbar t-C\hbar^{1/2}t^2\Bigr],
\]
and the bracket is at least $\tfrac{\sqrt\pi}4|S_0(\hbar)|$ for $t\ge t_*$, $t\le\hbar^{-1/5}$ and $\hbar\le\hbar_0$. Lemma~\ref{cl:amp}(iii) bounds $|S_0|$ below.
\end{proof}

